% !TEX encoding = UTF-8 Unicode
% !TeX program = xelatex
% !TEX TS-program = XeLaTeX

\documentclass[12pt, a4paper]{article}
%\documentclass[12pt, a5paper]{article}
\usepackage{mathtools}
\usepackage{unicode-math}
\usepackage{comment}
\usepackage{tikz}
\usepackage[pass]{geometry}
%\geometry{legalpaper, landscape, margin=1in}
\usepackage[empty]{fullpage}
\usepackage{setspace}
\usepackage[localise]{xepersian}
\setstretch{1.255}
\settextfont[Scale=1]{XB Niloofar}
\setlatintextfont[Scale=1]{XB Niloofar}%{Junicode} %Junicode
\setdigitfont{XB Niloofar}%{XB Zar}

\setlength{\parindent}{0ex}

\definecolor{answercolor}{HTML}{0080ff}
%\newcommand{\پاسخ}[2]{#1}
\newcommand{\پاسخ}[2]{\color{answercolor}\textbf{پاسخ: }#2\color{black}\vspace{-0.3cm}}

\begin{document}
	
	\begin{tikzpicture}[remember picture,overlay]
	\node[anchor=north,yshift=-1.5cm] at (current page.north) {\rl{به نام خدا}};
	\end{tikzpicture}%

\vspace{-1cm}
{\centering
	\textbf{پاسخ کوئیز اول درس ساختمان داده‌ها و الگوریتم‌ها}\\
مدرس: مهدی صفرنژاد       - ۱۳ اسفند ۹۷\\
\raggedleft\small
%نام و نام خانوادگی: \\
}
\vspace{-0.5cm}
\پرش‌متوسط\پرخط‌ا\\
%تمام سؤال‌ها جواب کوتاه هستند و توضیح یک‌خطی می‌خواهند.
\vspace{-1cm}
\پاراگراف{سؤال ۱. }
مقیاس‌پذیری یک الگوریتم به چه معنی است؟

\پاسخ{\vspace{0.5cm}}{
یعنی با بزرگ شدن اندازۀ ورودی، الگوریتم همچنان قابل استفاده باشد.

}

\پاراگراف{سؤال ۲.}
تست زمان را به صورت کوتاه تعریف کنید و ارتباط آن را با مطالب درس بیان کنید.

\پاسخ{\vspace{0.5cm}}{
یعنی مفید ماندن یک موضوع در طی زمان. مطالب این درس در چند دهۀ اخیر تست زمان را با موفقیت گذرانده‌اند.
	
}

\پاراگراف{سؤال ۳.}
این جمله صحیح است یا غلط؟ زمان اجرای یک الگوریتم برابر است با تعداد مراحل آن الگوریتم در مدل محاسباتی‌اش.

\پاسخ{\vspace{0.7cm}}{
غلط است. در برخی مدل‌های محاسبات مثل مدل پایتون یک مرحله ممکن است زمان زیادی صرف کند.
	
}

\پاراگراف{سؤال ۴.}
نام دو مدل محاسبات که در کلاس توضیح داده شدند را بنویسید.

\پاسخ{\vspace{0.5cm}}{
دو مورد از (۱) ماشین دسترسی تصادفی، (۲) ماشین اشاره‌گر و (۳) ماشین پایتون.
	
}

\پاراگراف{سؤال ۵.}
اگر $f(n) = O(g(n))$ باشد، آیا می‌توان استدلال کرد که
$2^{f(n)} = O(2^{g(n)})$؟

\پاسخ{\vspace{0.5cm}}{
خیر. اگر
$g(n)=\log n$
 و
$f(n)=2\log n=O(g(n))$
  باشد، داریم
$2^{f(n)}=n^2$ 
   و
$2^{g(n)}=n$.
     پس
$f(n)\neq O(2^{g(n)})$.
	
}

\پاراگراف{سؤال ۶.}
کدام‌یک از توابع زیر (به صورت مجانبی) بزرگ‌تر است؟
\vspace{-0.3cm}
$$f_1(n) = 2^{\log^4 n}, f_2(n) = n^{5\log^2 n}$$

\پاسخ{\vspace{0.5cm}}{
\vspace{-0.5cm}
	مبنای $\log$ها همان‌طور که سر کلاس گفته شد ۲ است.
$$f_1(n) = (2^{\log n})^{\log^3 n} = n^{\log^3 n} > n^{5\log^2 n} = f_2(n)$$
	\vspace{-1cm}
}

\پاراگراف{سؤال ۷.}
پیچیدگی زمانی تابع $T(n)$ با رابطۀ بازگشتی زیر چیست؟
\vspace{-0.3cm}
$$T(n)=T(\sqrt{n})+\log_2 n$$

\پاسخ{\vspace{0.5cm}}{
	\vspace{-1cm}
\begin{align*}
\begin{split}
T(n)&=\log_2 n + \log_2 (n^{1/2}) + \log_2 {n^{1/4}} + \dots \\
&= \log_2 n + (1/2)\log_2 n + (1/4)\log_2 n + \dots \leq 2\log_2 n = O(\log n)
\end{split}
\end{align*}

}

\پاراگراف{سؤال ۸.}
در مقایسۀ مرتب‌سازی درجی و ادغامی، یک مورد از برتری هر کدام را نام ببرید.

\پاسخ{\vspace{1cm}}{
مرتب‌سازی ادغامی زمان اجرای کمتری در بدترین حالت دارد. مرتب‌سازی درجی آرایه‌های تقریباً مرتب را سریع‌تر مرتب می‌کند. یا مرتب‌سازی درجی از حافظه اضافه استفاده نمی‌کند.
	
}

\پاراگراف{سؤال ۹.}
نام مرتب‌سازی مورد استفاده در زبان پایتون (الگوریتم تابع sort پایتون) چیست؟

\پاسخ{\vspace{0.5cm}}{
مرتب‌سازی تیم یا \متن‌لاتین{TimSort}.

}

\پاراگراف{سؤال ۱۰.}
آيا داده ساختار آرایۀ پویا برای واسط دنبالۀ پویا مناسب است؟

\پاسخ{\vspace{1.2cm}}{
خیر، زیرا درج در وسط دنباله به خاطر shift دادن زمان‌بر است.
	
}

\begin{flushleft}
موفق باشید
\end{flushleft}

\end{document}